\lecture{8}{18. September 2025}{Structural Analysis -- Indeterminate trusses}

\subsection{The method of sections}
When we are just interested in the force in a few members of the truss it is often found using the method of sections. This relies on the fact that for the truss to be in equilibrium then any part of the truss must also be in equilibrium. In this way, one can ``cut'' through a truss and analyze any part of it simply by applying the correct forces -- keep in mind that since we only have three equations of equilibrium then the cut should not intersect more than three members in which the forces are unknown. 

It is preferable to choose a point to sum the moments about which reduces the amount of forces the most possible. 
